\section{原子与光场的相互作用}
\label{ch:time-dependent-perturbation}

原子能级告诉我们哪些定态可以存在，光谱实验却测量态怎样在含时电磁场中改变。精确传播子包含所有时间排序；弱场只保留最低阶跃迁振幅，长观测时间把有限宽度的频率窗压成黄金法则，近共振窄带驱动则留下可相干往返的两个能级。再把辐射衰减和退相干加入，才得到光学 Bloch 方程。黄金法则、Rabi 振荡与 Bloch 动力学描述的是同一含时问题在不同时间尺度上的极限。
\subsection{形式精确传播、Magnus 展开与投影动力学}

设 Hilbert 空间上的自伴哈密顿量写成
\begin{equation}
H(t)=H_0+V(t),
\label{eq:tdpt-hamiltonian-split}
\end{equation}
其中 $H_0$ 为可解的参考哈密顿量，$V(t)$ 为含时相互作用。相互作用表象定义
\[
|\psi_I(t)\rangle=\ee^{\ii H_0(t-t_0)/\hbar}|\psi(t)\rangle,
\qquad
V_I(t)=\ee^{\ii H_0(t-t_0)/\hbar}V(t)\ee^{-\ii H_0(t-t_0)/\hbar}.
\]
传播子满足 Volterra 方程
\[
U_I(t,t_0)=I-\frac{\ii}{\hbar}\int_{t_0}^{t}V_I(t_1)U_I(t_1,t_0)\,\dd t_1,
\]
其形式精确解是时间序指数
\begin{equation}
U_I(t,t_0)=\mathcal T\exp\!\left[-\frac{\ii}{\hbar}\int_{t_0}^{t}V_I(t')\,\dd t'\right].
\label{eq:interaction-propagator}
\end{equation}
把 $U_I$ 迭代代回右端得到 Dyson 级数；若希望有限阶近似仍保持幺正性，则更适合写成 Magnus 形式 $U_I=\exp\Omega$。

在 $H_0|n\rangle=E_n|n\rangle$ 的本征基中，令
\[
|\psi(t)\rangle=\sum_n c_n(t)\ee^{-\ii E_n(t-t_0)/\hbar}|n\rangle,
\]
则没有任何微扰近似的系数方程为
\begin{equation}
\ii\hbar\dot c_m(t)
=\sum_n V_{mn}(t)\ee^{\ii\omega_{mn}(t-t_0)}c_n(t),
\qquad
\omega_{mn}=\frac{E_m-E_n}{\hbar}.
\label{eq:tdpt-exact-coefficient}
\end{equation}
后面的“一级微扰”“两能级截断”都必须说明是怎样从这组精确耦合方程中舍弃通道。

更一般地，若只关心一个有限维相关子空间，取正交投影 $P$ 与 $Q=I-P$。对 Schr\"odinger 方程作 $P/Q$ 分块并形式消去 $Q|\psi\rangle$，得到带记忆核的精确方程
\begin{equation}
\begin{aligned}
\ii\hbar\partial_t P|\psi(t)\rangle
={}&PHP\,P|\psi(t)\rangle
+PHQ\,U_Q(t,t_0)Q|\psi(t_0)\rangle\\
&-\frac{\ii}{\hbar}\int_{t_0}^{t}
PHQ\,U_Q(t,t')QHP\,P|\psi(t')\rangle\,\dd t',
\end{aligned}
\label{eq:two-level-memory-kernel}
\end{equation}
其中 $U_Q$ 是 $QHQ$ 在 $Q$ 子空间中的传播子。通常的两能级模型实质上是在能隙、失谐与耦合尺度允许时，把这类记忆和虚跃迁修正系统地压缩进有限维有效哈密顿量。

\begin{proof}[解]
Dyson 级数。将 Volterra 方程 $U_I=I-\frac{\ii}{\hbar}\int V_I U_I$ 迭代，得到 Dyson 级数写到二阶：
\[
U_I=I+U_1+U_2+O(V^3),
\]
\[
U_1=-\frac{\ii}{\hbar}\int_{t_0}^{t}V_I(t_1)\,\dd t_1,
\qquad
U_2=-\frac{1}{\hbar^2}\int_{t_0}^{t}\dd t_1\int_{t_0}^{t_1}\dd t_2\,
V_I(t_1)V_I(t_2).
\]
$U_2$ 的时序积分区域是三角形 $t_0\le t_2\le t_1\le t$，这正是时间序算符 $\mathcal T$ 的几何含义。

Magnus 指数展开。令 $U_I=\exp\Omega$，$\Omega=\Omega_1+\Omega_2+\cdots$ 按相互作用阶数展开。用 Baker--Campbell--Hausdorff 公式展开指数：
\[
U_I=I+\Omega_1+\Omega_2+\frac12\Omega_1^2+O(V^3).
\]
逐阶比较得到 $\Omega_1=U_1$ 以及
\[
\Omega_2=U_2-\frac12U_1^2.
\]

对称化方形积分。把 $U_1^2$ 的积分区域从方形 $[t_0,t]^2$ 拆成两个三角形 $t_1>t_2$ 与 $t_2>t_1$：
\[
\frac12U_1^2
=-\frac{1}{2\hbar^2}\int_{t_0}^{t}\dd t_1\int_{t_0}^{t_1}\dd t_2
\bigl[V_I(t_1)V_I(t_2)+V_I(t_2)V_I(t_1)\bigr].
\]
代入 $\Omega_2=U_2-\frac12U_1^2$，对称的 anticommutator 部分与 $U_2$ 抵消，只剩下对易子：
\begin{equation}
\Omega_2(t)
=-\frac{1}{2\hbar^2}\int_{t_0}^{t}\dd t_1\int_{t_0}^{t_1}\dd t_2\,
[V_I(t_1),V_I(t_2)].
\label{eq:magnus-second-order}
\end{equation}

高阶结构与幺正性。三阶项 $\Omega_3$ 涉及嵌套对易子 $\bigl[[V_I(t_1),V_I(t_2)],V_I(t_3)\bigr]$ 与 $\bigl[V_I(t_1),[V_I(t_2),V_I(t_3)]\bigr]$ 的组合。Magnus 展开的关键性质是：截断到任意有限阶，$\Omega$ 仍是反厄米的（当 $V_I$ 厄米），故 $U_I=\exp\Omega$ 自动幺正。这正是 Magnus 形式优于 Dyson 级数之处——Dyson 截断到 $O(V^2)$ 给出 $I+U_1+U_2$，一般不幺正；而 $\exp(\Omega_1+\Omega_2)$ 在二阶精度内保持概率守恒。

若不同时刻的相互作用彼此对易，$[V_I(t_1),V_I(t_2)]=0$，所有高阶 Magnus 对易子同时消失，时间序指数退化为普通指数
\[
U_I(t,t_0)=\exp\!\left[-\frac{\ii}{\hbar}\int_{t_0}^{t}V_I(t')\,\dd t'\right].
\]
这也给出脉冲面积定理的代数条件：若共振二能级系统的 $V_I(t)$ 始终正比于同一个 Pauli 矩阵，则跃迁概率只依赖 $\Theta=\int\Omega(t)\dd t$。若驱动方向随时间改变，二阶对易子会产生额外旋转。一个常用的充分收敛条件是 $\hbar^{-1}\int_0^t\|V_I(t')\|\,\dd t'<\pi$；长时间演化可分段应用 Magnus 展开。
\end{proof}

\subsection{弱场极限}

当跃迁概率在观测时间内远小于一，可在\eqnrefs{eq:tdpt-exact-coefficient}中把初态振幅取 $c_i\simeq1$、其余振幅取零阶值。于是从 $|i\rangle$ 到 $|f\rangle$ 的一级振幅为
\begin{equation}
c_f^{(1)}(t)
=-\frac{\ii}{\hbar}\int_{t_0}^{t}
V_{fi}(t')\ee^{\ii\omega_{fi}(t'-t_0)}\,\dd t'.
\label{eq:first-order-transition-amplitude}
\end{equation}
它首先是外场在 Bohr 频率处的有限时间 Fourier 分量，而不是一个“能量守恒 \(\delta\) 函数”。只有连续谱、长观测时间并且矩阵元在共振窗口内足够平滑时，才出现 Fermi 黄金法则
\begin{equation}
\Gamma_{i\to f}
=\frac{2\pi}{\hbar}|V_{fi}|^2\rho(E_f),
\label{eq:fermi-golden-rule}
\end{equation}
其中 $\rho(E_f)$ 是按能量计的末态态密度。黄金法则的控制条件包括弱跃迁、连续或准连续末态、关联时间远短于总体衰减时间；在离散强耦合体系中它会被相干 Rabi 交换取代。

\begin{proof}[解]
一阶微扰振幅。取在 $0<t<T$ 内近似恒定的矩阵元 $V_{fi}$。由\eqnrefs{eq:first-order-transition-amplitude}，对初态 $|i\rangle$ 的一阶跃迁振幅
\begin{equation*}
 c_f^{(1)}(T)
 =-\frac{\ii V_{fi}}{\hbar}\int_0^T\ee^{\ii\omega_{fi}t}\,\dd t
 =-\frac{\ii V_{fi}}{\hbar}
 \ee^{\ii\omega_{fi}T/2}
 \frac{2\sin(\omega_{fi}T/2)}{\omega_{fi}},
\end{equation*}
其中 $\omega_{fi}=(E_f-E_i)/\hbar$ 是 Bohr 频率。

跃迁概率。取模平方，
\begin{equation*}
 P_{i\to f}(T)
 =|c_f^{(1)}(T)|^2
 =\frac{|V_{fi}|^2}{\hbar^2}
 \frac{4\sin^2(\omega_{fi}T/2)}{\omega_{fi}^2}.
\end{equation*}
被积函数是 sinc 平方型共振峰，中心在 $\omega_{fi}=0$（能量守恒），宽度 $\sim1/T$。

连续末态求和。对连续末态求和，用 $E_f-E_i=\hbar\omega_{fi}$ 换元：
\begin{equation*}
 \frac{P(T)}{T}
 =\int\dd E_f\,\rho(E_f)
 \frac{|V_{fi}|^2}{\hbar^2}
 \frac{4\sin^2[(E_f-E_i)T/(2\hbar)]}{(E_f-E_i)^2T/\hbar^2}.
\end{equation*}

分布极限。关键恒等式是 sinc 平方函数的分布极限：
\begin{equation*}
 \lim_{T\to\infty}
 \frac{4\sin^2(\omega T/2)}{\omega^2T}
 =2\pi\delta(\omega).
\end{equation*}
验证：对测试函数 $\varphi(\omega)$，$\int\frac{4\sin^2(\omega T/2)}{\omega^2T}\varphi(\omega)\dd\omega\to2\pi\varphi(0)$（换元 $u=\omega T/2$ 后 Riemann--Lebesgue）。代入第三步，
\begin{equation*}
 \Gamma_{i\to f}
 =\lim_{T\to\infty}\frac{P(T)}{T}
 =\int\dd E_f\,\rho(E_f)|V_{fi}|^2\frac{2\pi}{\hbar}\delta(E_f-E_i)
 =\frac{2\pi}{\hbar}|V_{fi}|^2\rho(E_i),
\end{equation*}
即\eqnrefs{eq:fermi-golden-rule}。

有限 $T$ 时，共振峰宽度量级为 $1/T$，所以真正被近似为常数的是这一窄频窗内的 $V_{fi}$ 与 $\rho$。黄金法则要求弱跃迁和连续或准连续末态；离散强耦合体系中，它被相干 Rabi 交换取代。


\end{proof}

对与黑体辐射场弱耦合的两能级原子，若 $u_\nu$ 是“每单位普通频率”的谱能量密度，受激吸收、受激辐射和自发辐射率分别写为 $B_{12}u_\nu$、$B_{21}u_\nu$、$A_{21}$。热平衡与 Planck 分布给出
\begin{equation}
 g_1B_{12}=g_2B_{21},
\qquad
A_{21}=\frac{8\pi h\nu^3}{c^3}B_{21}.
\label{eq:einstein-relations}
\end{equation}
若改用每单位角频率的谱密度 $u_\omega$，由于 $u_\omega\dd\omega=u_\nu\dd\nu$，相应 $B$ 系数的数值定义随谱密度约定改变；物理跃迁率保持不变。

在长波电偶极近似中
\[
V_{E1}(t)=-\mathbf d\cdot\mathbf E(t),
\qquad
\mathbf d=\sum_a q_a\mathbf r_a.
\]
球分量 $d_q$ 是奇宇称秩-$1$ 张量，因此\eqnrefs{eq:wigner-eckart}立即给出
\begin{equation}
\Delta M=q=0,\pm1,
\qquad
\Delta J=0,\pm1,
\qquad
J=0\not\leftrightarrow J'=0,
\label{eq:e1-angular-selection}
\end{equation}
并要求初末态宇称相反。对单电子中心势本征态还可进一步写成 $\Delta l=\pm1$；在良好的 LS 耦合且电偶极算符不作用于自旋时有 $\Delta S=0$。这些是额外结构下的特例，不应与一般总角动量选择定则混为一谈。

\begin{proof}[解]
由 Wigner--Eckart 定理，秩-1 张量 $d_q$ 的矩阵元
\begin{equation*}
 \langle \alpha'J'M'|d_q|\alpha JM\rangle
 =(-1)^{J'-M'}
 \begin{pmatrix}J'&1&J\\-M'&q&M\end{pmatrix}
 \langle\alpha'J'||d||\alpha J\rangle.
\end{equation*}
3-j 符号非零有两个必要条件。其一为磁量子数守恒：三行第二列之和为零，
\begin{equation*}
 -M'+q+M=0
 \quad\Longrightarrow\quad
 M'=M+q,
 \qquad q=0,\pm1,
\end{equation*}
即 $\Delta M=0,\pm1$。其二为三角条件：$(J',1,J)$ 可构成三角形，
\begin{equation*}
 |J-1|\le J'\le J+1,
 \qquad
 J'\in\{J-1,J,J+1\},
\end{equation*}
即 $\Delta J=0,\pm1$。若 $J=J'=0$，三角条件要求 $|0-1|\le0\le0+1$，左端 $1\le0$ 失败，因此 $0\leftrightarrow0$ 被排除。

另一方面空间反演算符 $\Pi$ 满足 $\Pi^2=I$ 且 $\Pi\mathbf d\Pi^{-1}=-\mathbf d$，因为电偶极矩是极矢量。若初末态是宇称本征态，
\begin{equation*}
 \Pi|i\rangle=\pi_i|i\rangle,
 \qquad
 \Pi|f\rangle=\pi_f|f\rangle,
 \qquad
 \pi_i,\pi_f\in\{+1,-1\},
\end{equation*}
则
\begin{align*}
 \langle f|\mathbf d|i\rangle
 &=\pi_f\pi_i\,\langle f|\Pi^{-1}\mathbf d\Pi|i\rangle\\
 &=\pi_f\pi_i\,\langle f|(-\mathbf d)|i\rangle\\
 &=-\pi_f\pi_i\,\langle f|\mathbf d|i\rangle.
\end{align*}
移项得
\begin{equation*}
 (1+\pi_f\pi_i)\langle f|\mathbf d|i\rangle=0,
\end{equation*}
非零矩阵元要求 $1+\pi_f\pi_i=0$，即 $\pi_f\pi_i=-1$，初末态宇称相反。角动量约束来自旋转对称性，宇称约束来自空间反演对称性，两者是独立的对称性操作，必须同时满足；仅满足其一的跃迁仍被禁戒。
\end{proof}

例：氢原子 $1s\to2p$ 初末态宇称相反，电偶极跃迁允许；$1s\to2s$ 宇称相同，电偶极跃迁禁戒。钠 $D$ 线 $3s\to3p$ 同时满足角动量与宇称选择定则。

\subsection{近共振截断、旋波近似与 Rabi 动力学}

设一个孤立近共振双态子空间由基 $\{|e\rangle,|g\rangle\}$ 张成，裸跃迁频率为 $\omega_0=(E_e-E_g)/\hbar$。所谓“两能级近似”要求外场带宽和 Rabi 频率远小于到其他能级的失谐，并且消去的能级只产生可控的虚跃迁修正。取实电场的复包络约定
\begin{equation}
\mathbf E(t)
=\frac12\left[\boldsymbol{\mathcal E}(t)\ee^{-\ii\omega_L t}
+\boldsymbol{\mathcal E}^{*}(t)\ee^{\ii\omega_L t}\right],
\qquad
\Omega(t)=-\frac{\mathbf d_{eg}\cdot\boldsymbol{\mathcal E}(t)}{\hbar},
\label{eq:rabi-frequency-convention}
\end{equation}
并定义失谐 $\Delta=\omega_L-\omega_0$。在旋转表象中，保留缓慢项而舍弃频率约为 $\omega_L+\omega_0$ 的反旋项，需要
\begin{equation}
|\Omega|,|\Delta|,\tau_{\rm env}^{-1}\ll \omega_L+\omega_0,
\label{eq:rwa-condition}
\end{equation}
其中 $\tau_{\rm env}$ 是包络变化时间尺度。到该阶的旋转波近似哈密顿量可写为
\begin{equation}
H_{\rm RWA}
=\frac{\hbar}{2}
\begin{pmatrix}
-\Delta&\Omega\\
\Omega^*&\Delta
\end{pmatrix},
\label{eq:rwa-hamiltonian}
\end{equation}
忽略与单位阵成正比的公共能量。

当 $\Omega$ 和 $\Delta$ 为常数时，广义 Rabi 频率
\[
\Omega_R=\sqrt{|\Omega|^2+\Delta^2}
\]
给出从 $|g\rangle$ 出发的激发概率
\begin{equation}
P_e(t)
=\frac{|\Omega|^2}{\Omega_R^2}
\sin^2\!\left(\frac{\Omega_Rt}{2}\right).
\label{eq:rabi-probability}
\end{equation}
共振且驱动相位固定时，若不同时刻哈密顿量彼此对易，则动力学只依赖 pulse area $\Theta=\int|\Omega(t)|\dd t$；一般失谐或时变相位下不能把“脉冲面积”当成普适定理。

远失谐 $|\Delta|\gg|\Omega|$ 时，可对\eqnrefs{eq:rwa-hamiltonian}作二阶非简并微扰，得到 AC Stark 位移
\begin{equation}
\delta E_g=+\frac{\hbar|\Omega|^2}{4\Delta},
\qquad
\delta E_e=-\frac{\hbar|\Omega|^2}{4\Delta},
\label{eq:ac-stark-shift}
\end{equation}
这里的正负号与\eqnrefs{eq:rabi-frequency-convention}和 $\Delta=\omega_L-\omega_0$ 的约定绑定；改用相反失谐定义时必须整体同步改变。

\begin{proof}[解]
实验室表象哈密顿量。在两能级子空间取 $\sigma_z=|e\rangle\langle e|-|g\rangle\langle g|$、$\sigma_+=|e\rangle\langle g|$、$\sigma_-=|g\rangle\langle e|$。去掉公共能量 $(E_e+E_g)/2$ 后，实验室表象哈密顿量为
\begin{equation*}
 H=\frac{\hbar\omega_0}{2}\sigma_z
 +\frac{\hbar}{2}
 \left[\Omega\ee^{-\ii\omega_L t}+\Omega^*\ee^{\ii\omega_L t}\right]
 \frac{\sigma_++\sigma_-}{2}.
\end{equation*}
展开得共振项 $\Omega\sigma_\pm\ee^{\mp\ii\omega_Lt}/2$ 与反旋项 $\Omega\sigma_\mp\ee^{\mp\ii\omega_Lt}/2$。

用 $R(t)=\exp(+\ii\omega_Lt\sigma_z/2)$ 变到以激光频率旋转的表象，
\begin{equation*}
 H_{\rm rot}=RHR^\dagger+\ii\hbar\dot R R^\dagger.
\end{equation*}
$R(t)$ 把 $\sigma_\pm$ 旋转为 $\sigma_\pm\ee^{\mp\ii\omega_Lt}$。于是耦合项中，共振部分（$\Omega\ee^{-\ii\omega_Lt}\cdot\sigma_+\ee^{+\ii\omega_Lt}=\Omega\sigma_+$）变为常数，反旋部分携带 $\ee^{\pm2\ii\omega_Lt}$ 的快振荡。$\ii\hbar\dot RR^\dagger=-\hbar\omega_L\sigma_z/2$ 项将裸能隙 $\omega_0$ 替换为失谐 $\Delta=\omega_0-\omega_L$。在\eqnrefs{eq:rwa-condition}（$\omega_L\gg\Omega$）下丢弃快振荡项，得到\eqnrefs{eq:rwa-hamiltonian}；修正量级为 $(\Omega/\omega_L)^2$，对应 Bloch--Siegert 移位 $\delta_{\rm BS}=\Omega^2/[4(\omega_L+\omega_0)]$。

对常系数旋转波近似哈密顿量
\begin{equation*}
 H_{\rm RWA}=\frac{\hbar}{2}\bigl(\Delta\,\sigma_z+\Omega\,\sigma_++\Omega^*\,\sigma_-\bigr)
\end{equation*}
满足
\begin{equation*}
 H_{\rm RWA}^2=\frac{\hbar^2\Omega_R^2}{4}I,
 \qquad
 \Omega_R=\sqrt{|\Omega|^2+\Delta^2},
\end{equation*}
（Pauli 矩阵的 Clifford 代数关系 $\sigma_i\sigma_j=\delta_{ij}I+\ii\epsilon_{ijk}\sigma_k$）。因此指数截断为有限和：
\begin{equation*}
 U(t)=\exp\!\left(-\frac{\ii t}{\hbar}H_{\rm RWA}\right)
 =\cos\frac{\Omega_Rt}{2}\,I
 -\frac{2\ii}{\hbar\Omega_R}
 \sin\frac{\Omega_Rt}{2}\,H_{\rm RWA}.
\end{equation*}

令初态为 $|g\rangle$，取 $\langle e|U(t)|g\rangle$ 的模平方。由 $\langle e|\sigma_+|g\rangle=1$、$\langle e|\sigma_z|g\rangle=0$，
\begin{equation*}
 P_e(t)=|\langle e|U(t)|g\rangle|^2
 =\frac{|\Omega|^2}{\Omega_R^2}\sin^2\frac{\Omega_R t}{2},
\end{equation*}
即\eqnrefs{eq:rabi-probability}。共振时 $\Delta=0$ 布居数完全翻转，远失谐时翻转幅度 $\sim|\Omega/\Delta|^2$，主要效应为 AC Stark 移位。
\end{proof}

\subsection{开放驱动两能级系统}

真实原子与真空电磁场、碰撞环境和激光相位噪声耦合。若环境关联时间远短于系统弛豫时间，并采用适当 Born--马尔可夫与久期近似，约化态满足 GKSL 方程。对自发辐射率 $\Gamma_1$ 和纯退相干率 $\gamma_\phi$，
\begin{equation}
\dot\rho
=-\frac{\ii}{\hbar}[H_{\rm RWA},\rho]
+\Gamma_1\mathcal D[\sigma_-]\rho
+\frac{\gamma_\phi}{2}(\sigma_z\rho\sigma_z-\rho),
\label{eq:optical-bloch-master}
\end{equation}
其中 $\mathcal D[L]\rho=L\rho L^\dagger-\tfrac12\{L^\dagger L,\rho\}$，横向退相干率为
\[
\Gamma_2=\frac{\Gamma_1}{2}+\gamma_\phi.
\]
写 $u=2\Re\rho_{eg}$、$v=2\Im\rho_{eg}$、$w=\rho_{ee}-\rho_{gg}$，并通过重定义基态相位把常数 $\Omega$ 取为实数，可得光学 Bloch 方程
\begin{equation}
\dot u=-\Gamma_2u+\Delta v,
\qquad
\dot v=-\Gamma_2v-\Delta u-\Omega w,
\qquad
\dot w=-\Gamma_1(w+1)+\Omega v.
\label{eq:optical-bloch-equations}
\end{equation}
稳态激发占据为
\begin{equation}
\rho_{ee}^{\rm ss}
=\frac{|\Omega|^2\Gamma_2}
{2\left[\Gamma_1(\Gamma_2^2+\Delta^2)+|\Omega|^2\Gamma_2\right]}.
\label{eq:bloch-steady-excited}
\end{equation}
若只有自发辐射，$\Gamma_2=\Gamma_1/2\equiv\Gamma/2$，则
\begin{equation}
\rho_{ee}^{\rm ss}
=\frac{|\Omega|^2}{\Gamma^2+4\Delta^2+2|\Omega|^2}
=\frac{s_0/2}{1+s_0+(2\Delta/\Gamma)^2},
\qquad
s_0=\frac{2|\Omega|^2}{\Gamma^2}.
\label{eq:two-level-saturation}
\end{equation}
下一章以后用到的散射力、Doppler 冷却和精密操控都以这一稳态响应作为内部态输入；若饱和很强、环境非马尔可夫或存在多重简并暗态，则必须回到更完整的主方程。

\begin{proof}[解]
令\eqnrefs{eq:optical-bloch-equations} 左端为零，稳态满足
\begin{align*}
 0&=\Delta v-\Gamma_2 u,\\
 0&=-\Delta u-\Gamma_2 v-\Omega w,\\
 0&=-\Gamma_1(w+1)+\Omega v.
\end{align*}
由第一式解出
\begin{equation*}
 u=\frac{\Delta}{\Gamma_2}v.
\end{equation*}
代回第二式，
\begin{equation*}
 0=-\Delta\cdot\frac{\Delta}{\Gamma_2}v-\Gamma_2 v-\Omega w
 =-\left(\frac{\Delta^2}{\Gamma_2}+\Gamma_2\right)v-\Omega w,
\end{equation*}
整理得
\begin{equation*}
 0=-\frac{\Gamma_2^2+\Delta^2}{\Gamma_2}v-\Omega w,
\end{equation*}
于是
\begin{equation*}
 v=-\frac{\Omega\Gamma_2}{\Gamma_2^2+\Delta^2}w.
\end{equation*}
再代入第三式，
\begin{equation*}
 \Gamma_1(w+1)
 =\Omega v
 =-\frac{\Omega^2\Gamma_2}{\Gamma_2^2+\Delta^2}w.
\end{equation*}
把含 $w$ 的项合并到一端，
\begin{equation*}
 \Gamma_1 w+\frac{\Omega^2\Gamma_2}{\Gamma_2^2+\Delta^2}w
 =-\Gamma_1,
\end{equation*}
解得
\begin{equation*}
 w=-\frac{\Gamma_1}{\Gamma_1+\dfrac{\Omega^2\Gamma_2}{\Gamma_2^2+\Delta^2}}
 =-\frac{\Gamma_1(\Gamma_2^2+\Delta^2)}
 {\Gamma_1(\Gamma_2^2+\Delta^2)+\Omega^2\Gamma_2}.
\end{equation*}
由 $\rho_{ee}=(1+w)/2$，代入并化简即得\eqnrefs{eq:bloch-steady-excited}。当 $|\Omega|\to\infty$ 时，分母中 $\Omega^2\Gamma_2$ 主导，故 $w\to0$，从而 $\rho_{ee}\to1/2$，这不是"激发失败"，而是受驱二能级系统在强饱和下上下能级占据趋于相等。
\end{proof}

\subsection{典型原子--光系统}
\label{sec:rabi-bloch-experiment}

黄金法则、Rabi 频率和 Bloch 方程中的速率只有代入具体原子的线宽、偶极矩和激光强度后才可比较。下面用几类常见原子--光系统估计这些量，并指出弱场、近共振和稳态条件各自覆盖的实验范围。

二能级系统的关键参数为 Rabi 频率 $\Omega=\vb d\cdot\vb E/\hbar$、自发辐射率 $\Gamma$ 和失谐 $\Delta$。对常见原子：
\begin{center}
\begin{tabular}{cccccc}
\hline
原子 & 跃迁 & $\lambda$ (nm) & $\Gamma/2\pi$ (MHz) & $d$ (ea$_0$) & $\tau$ (ns) \\
\hline
H & $1s\to2p$ & 121.6 & 99.6 & 0.74 & 1.60 \\
Na & $3s\to3p$ & 589.0 & 6.06 & 0.71 & 26.2 \\
Rb & $5s\to5p$ & 780.2 & 6.07 & 0.68 & 26.2 \\
Cs & $6s\to6p$ & 852.3 & 5.22 & 0.69 & 30.5 \\
Ca & $4s^2\to4s4p$ & 422.7 & 37.4 & 1.49 & 4.25 \\
\hline
\end{tabular}
\end{center}
其中 $a_0=0.529\times10^{-10}\,\mathrm{m}$ 为 Bohr 半径。氢 $2p\to1s$ 的 Einstein $A$ 系数 $A_{21}=\Gamma=6.27\times10^8\,\mathrm{s}^{-1}$，对应寿命 $\tau=1.60\,\mathrm{ns}$。

共振驱动 ($\Delta=0$) 下二能级系统的激发态布居
\[
 P_e(t)=\sin^2\frac{\Omega t}{2}
\]
在 $\Gamma=0$ 时为完美振荡。有阻尼时 ($\Gamma>0$)，振荡以速率 $\Gamma/2$ 衰减，稳态趋于饱和值。典型实验参数：
\begin{itemize}
 \item Rb 原子在 $I=1\,\mathrm{mW/cm^2}$ 激光下：$\Omega/2\pi\approx5\,\mathrm{MHz}$，$\Omega/\Gamma\approx0.8$（中等饱和）。
 \item 强场 $I=10\,\mathrm{W/cm^2}$：$\Omega/2\pi\approx500\,\mathrm{MHz}\gg\Gamma$，$\pi$ 脉冲时间 $t_\pi=\pi/\Omega\approx1\,\mathrm{ns}$。
 \item 冷原子实验中已观测到上千次相干 Rabi 振荡（退相干时间 $T_2\sim100\,\mathrm{\mu s}$）。
\end{itemize}

弱驱动 ($\Omega\ll\Gamma$) 下 Fermi 黄金法则给出跃迁率 $W=\Omega^2\Gamma/(\Gamma^2+\Delta^2)$，激发态布居 $P_e\approx Wt\propto t$ 线性增长。强驱动 ($\Omega\gtrsim\Gamma$) 时 Rabi 振荡显著，黄金法则失效——需用完整 Bloch 方程。判据：黄金法则要求 $\Omega^2/(\Gamma^2+\Delta^2)\ll1$（饱和参数 $s\ll1$）。

原子蒸气在连续激光扫描下，强驱动使 $\rho_{ee}\to1/2$，吸收系数 $\alpha\propto(1-2\rho_{ee})$ 下降。在共振频率处饱和最强、吸收最弱，形成饱和吸收峰的「反兰姆凹陷」。钠 D 线饱和吸收光谱的典型线宽 $\sim50\,\mathrm{MHz}$（功率加宽），用于激光频率锁定至原子跃迁——这是原子钟和激光冷却的标准技术。

Bloch 矢量 $\vb R=(u,v,w)$（$u=\rho_{ge}+\rho_{eg}$, $v=-i(\rho_{ge}-\rho_{eg})$, $w=\rho_{ee}-\rho_{gg}$）在 Bloch 球上运动。无阻尼时绕有效场 $\vb\Omega_{\rm eff}=(\Omega,0,\Delta)$ 做刚体旋转；有阻尼时 Bloch 矢量收缩至稳态。几何图像使量子态演化可视化：$\pi$ 脉冲把 $\vb R$ 从北极翻到南极（完全激发），$\pi/2$ 脉冲移至赤道（叠加态），自由演化绕 $z$ 轴以频率 $\Delta$ 进动。这一图像是 NMR 和量子比特操控的标准语言。

\begin{proof}[解]
二能级哈密顿量
\begin{equation*}
  H(t)=\frac{1}{2}\begin{pmatrix}\alpha t&\Delta\\\Delta&-\alpha t\end{pmatrix},
\end{equation*}
能级差 \(\epsilon(t)=\alpha t\) 线性穿越零点，耦合 \(\Delta\) 恒定；绝热性由 \(\Delta^2/(\hbar\alpha)\) 控制。Schrödinger 方程化为 Weber 方程（抛物柱函数）。设 \(t\to-\infty\) 处于下能级 \(\psi_-\)，\(t\to+\infty\) 时处于上能级的概率
\begin{equation}
  P_{\rm LZ}=\exp\!\left(-\frac{\pi\Delta^2}{2\hbar\alpha}\right).
  \label{eq:tdpt-landau-zener}
\end{equation}
证明将方程化为 Weber 方程后由 \(t\to\pm\infty\) 连接公式读 Stokes 系数即得。绝热极限 \(\alpha\to0\) 时 \(P_{\rm LZ}\to0\)，突然极限 \(\alpha\to\infty\) 时 \(P_{\rm LZ}\to1\)。多次穿越时各振幅连同动力学相位相干叠加，形成 Stückelberg 干涉。

周期驱动由 Floquet 定理刻画：\(U(t,0)=P(t)\ee^{-iH_F t/\hbar}\)，\(P(t+T)=P(t)\) 为周期算符，准能 \(E_F\) 模 \(\hbar\omega\) 定义。在扩展空间 \(\mathcal H\otimes\mathcal T\) 上 \(H_F=H(t)-i\hbar\partial_t\)，将含时问题化为时间无关本征值问题。
\end{proof}
